\documentclass[11pt]{article}
\usepackage[T1]{fontenc}
\usepackage[margin=1in]{geometry}
\usepackage{amsmath,amssymb}
\usepackage[hidelinks]{hyperref}

\title{Literature Status: The Lorentz \texttt{lorentzck} Conjecture}
\author{Run \texttt{fa\_banach\_001}, lane 11}
\date{2026-07-18}

\begin{document}
\maketitle

\section*{Status}

This is a lightweight literature-status packet.  The Lorentz-space conjecture
from Defant, Mastyl\l{}o, and Michels, arXiv:math/0006034, is already answered
in later s-number literature by Hinrichs and by Hinrichs--Michels.  This packet
is not an original proof claim by the run.

\section*{Source Conjecture}

The source paper is Andreas Defant, Mieczyslaw Mastyl\l{}o, and Carsten
Michels, \emph{Summing inclusion maps between symmetric sequence spaces},
arXiv:math/0006034; published in \emph{Transactions of the American
Mathematical Society} 354 (2002), 4473--4492.

Near the end of the paper, on page 19 of the locally rendered
\texttt{source\_paper.pdf} and after the corollary containing the displayed
formula labelled \texttt{lorentzck}, the authors conjecture that
\[
  c_k\!\left(\operatorname{id}:
    \ell_{v_1,v_2}^n \to \ell_{u_1,u_2}^n\right)
  \asymp
  a_k\!\left(\operatorname{id}:
    \ell_{v_1,v_2}^n \to \ell_{u_1,u_2}^n\right)
  \asymp
  (n-k+1)^{1/u_1-1/v_1}
\]
for every \(1<u_1<v_1<2\) and every \(1\leq u_2,v_2\leq\infty\).  The source
paper proves only special secondary-parameter ranges immediately before this
sentence.

\section*{Supporting Literature}

The approximation-number side is covered by Aicke Hinrichs,
\emph{Approximation Numbers of Identity Operators between Symmetric Sequence
Spaces}, \emph{Journal of Approximation Theory} 118 (2002), 305--315.  Its
abstract states that it proves asymptotic formulas for approximation quantities
of identity operators between symmetric sequence spaces and gives applications
to Lorentz and Orlicz sequence spaces, extending Defant--Mastyl\l{}o--Michels.

The Gelfand-number side is covered by Aicke Hinrichs and Carsten Michels,
\emph{Gelfand Numbers of Identity Operators Between Symmetric Sequence Spaces},
\emph{Positivity} 10 (2006), 111--133.  The introduction explicitly identifies
the Defant--Mastyl\l{}o--Michels conjecture as formula (1.1), involving both
\(a_k\) and \(c_k\), and states that the paper proves that formula under its
interpolation hypotheses.  The paper also notes that the displayed formulas may
be read with approximation numbers by replacing the Gelfand/Kolmogorov symbols
with \(a_k\).

For the present target, Example 7.5(ii) in the 2006 paper gives
\[
  c_k\!\left(\operatorname{id}:
    \ell_{q_1,q_2}^n \to \ell_{p_1,p_2}^n\right)
  \asymp
  (n-k+1)^{1/p_1-1/q_1}
\]
for \(0<p_1<q_1<\infty\) and \(0<p_2,q_2\leq\infty\).

\section*{Identification}

Set
\[
  p_1=u_1,\qquad q_1=v_1,\qquad p_2=u_2,\qquad q_2=v_2 .
\]
The supporting range \(0<p_1<q_1<\infty\), \(0<p_2,q_2\leq\infty\) contains the
source range \(1<u_1<v_1<2\), \(1\leq u_2,v_2\leq\infty\).  Under this
substitution, the exponent \(1/p_1-1/q_1\) is exactly \(1/u_1-1/v_1\), and the
identity map is precisely
\[
  \operatorname{id}:\ell_{v_1,v_2}^n\to\ell_{u_1,u_2}^n .
\]
Together with the approximation-number literature cited above, this answers
the full displayed \texttt{lorentzck} conjecture from the source paper.

\section*{Search Evidence and Scope}

Cheap run indexes showed no prior exact packet for arXiv:math/0006034.  A
bounded literature search for the source arXiv id, \texttt{lorentzck}, and the
Hinrichs/Hinrichs--Michels titles located the two decisive supporting papers.
The 2006 paper explicitly connects itself to the Defant--Mastyl\l{}o--Michels
conjecture, so this packet is classified as \texttt{literature\_already\_answered}.

No novelty is claimed.  The recorded result is only the already-known answer to
the Lorentz \texttt{lorentzck} conjecture; unrelated questions from the source
paper are outside this packet.

\section*{References}

\begin{enumerate}
\item A. Defant, M. Mastyl\l{}o, and C. Michels, \emph{Summing inclusion maps
between symmetric sequence spaces}, Trans. Amer. Math. Soc. 354 (2002),
4473--4492. DOI: \href{https://doi.org/10.1090/S0002-9947-02-03056-8}
{10.1090/S0002-9947-02-03056-8}.
\item A. Hinrichs, \emph{Approximation Numbers of Identity Operators between
Symmetric Sequence Spaces}, J. Approx. Theory 118 (2002), 305--315. DOI:
\href{https://doi.org/10.1006/jath.2002.3726}{10.1006/jath.2002.3726}.
\item A. Hinrichs and C. Michels, \emph{Gelfand Numbers of Identity Operators
Between Symmetric Sequence Spaces}, Positivity 10 (2006), 111--133. DOI:
\href{https://doi.org/10.1007/s11117-005-0022-1}
{10.1007/s11117-005-0022-1}.
\end{enumerate}

\end{document}
